\section{Report structure}

The remainder of this report is organized into four chapters.

\begin{itemize}
    \item \textbf{Chapter 2 (Background)} reviews the theoretical foundations underpinning this work: the role of Small Language Models in the legal domain, knowledge distillation and its exposure-bias limitation, parameter-efficient fine-tuning (LoRA and QLoRA), reinforcement-learning-based alignment and preference optimization, and the evaluation metrics used in legal artificial intelligence.
    \item \textbf{Chapter 3 (Methods and Models)} presents the proposed framework. It formalizes the legal-reasoning task, gives an overview of the two-track pipeline, and details each phase in turn: Phase 1 (Offline Logit Distillation), Phase 2 (On-Policy Knowledge Distillation), and Phase 3 (Diagnosis-Driven Preference Optimization).
    \item \textbf{Chapter 4 (Experiments)} describes the datasets, the experimental setup and hyperparameter configuration, and reports the empirical results, ablation studies, and qualitative analysis on both in-distribution and out-of-distribution benchmarks.
    \item \textbf{Chapter 5 (Conclusion)} summarizes the principal findings and contributions, discusses the limitations of the study, and outlines directions for future work.
\end{itemize}
